\begin{frame}{Normalisation des données}

	\begin{definitionbox}
		La \textbf{normalisation} est un processus de structuration des données dans une base de données relationnelle pour minimiser la redondance et éviter les anomalies de mise à jour, d'insertion et de suppression.
	\end{definitionbox}

	Elle consiste à diviser les données en plusieurs tables liées entre elles par des clés de jointure, en respectant des règles appelées \textit{formes normales}.

\end{frame}

\begin{frame}{Comment normaliser ?}

	En informatique, une \textbf{forme normale} est un ensemble de règles visant à organiser les données dans une base de données relationnelle afin de minimiser la redondance et d'améliorer l'intégrité des données.

	(faire diagramme de cercles concentriques 1FN, 2FN, 3FN, BCFM, 4FN, 5FN)

	\note{En pratique on s'arrête à la 3FN.}

\end{frame}

\begin{frame}{Première forme normale}

	\begin{definitionbox}
		Une relation est en \textbf{première forme normale} (1FN) si et seulement si ses attributs sont atomiques.
	\end{definitionbox}

	Un attribut est \textbf{atomique} s'il ne peut pas être décomposé sans perte d'information essentielle.

	\examplebox{
        Attributs atomiques : date, prix

		Attributs non atomiques : liste d'ingrédients
    }

	Un numéro de sécurité sociale ne peut pas être atomique.

\end{frame}

\begin{frame}{Deuxième forme normale}

	\begin{definitionbox}
		Une relation est en \textbf{deuxième forme normale} (2FN) si et seulement elle est en 1FN et tout attribut non clé est en dépendance fonctionnelle \underline{élémentaire} avec la clé.
	\end{definitionbox}

\end{frame}

\begin{frame}{Troisième forme normale}

	\begin{definitionbox}
		Une relation est en \textbf{troisième forme normale} (3FN) si et seulement elle est en 2FN et tout attribut non clé est en dépendance fonctionnelle \underline{directe} avec la clé.
	\end{definitionbox}

\end{frame}

\begin{frame}{Formes normales}

    % https://www.datacamp.com/tutorial/normalization-in-sql
    \begin{itemize}
        \item \textbf{1NF} : élimination des groupes de répétition
        \item \textbf{2NF} : élimination des dépendances partielles
        \item \textbf{3NF} : élimination des dépendances transitives
        \item \textbf{BCNF} : Boyce-Codd Normal Form
        \item \textbf{4NF} : élimination des dépendances multivaluées
        \item \textbf{5NF} : élimination des dépendances de jointure
%        \item \textbf{DNF} : forme dénormalisée
%        \item \textbf{ENF} : forme normalisée étendue
%        \item \textbf{ONF} : forme normale optimale
%        \item \textbf{TNF} : forme normale temporelle
%        \item \textbf{INF} : forme normale d'inclusion
%        \item \textbf{PNF} : forme normale de projection
    \end{itemize}

\end{frame}

\begin{frame}{Application de formes normales}

    Exemple de formes normales sur une table de base de données relationnelle.

\end{frame}

\begin{frame}{Comment normaliser la table de client$\cdot$es ?}

	\begin{table}[]
	\resizebox{\textwidth}{!}{
		\begin{tabular}{|c|c|c|c|c|c|c|}
			\hline
			\rowcolor{lipstick}\color{white}\textbf{Nom} & \color{white}\textbf{Âge} & \color{white}\textbf{Département} & \color{white}\textbf{Région} & \color{white}\textbf{Dernière visite}  \\\hline
			Alice Fleury & 30 & Paris & Île-de-France & 01/02/2023 \\
			Bob Brown & 39 & Var & CCA & 30/03/2026 \\
			Charlie O'Brien & 26 & Finistère & Bretagne & 25/06/2023 \\
			Daniel Doyle & 42 & Ille-et-Vilaine & Bretagne & 02/03/2022 \\
			Edna McCarthy & 58 & Oise & Hauts-de-France & 06/06/2025 \\
			Fabienne Lenoir & 66 & Savoie & Auvergne-Rhône-Alpes & 12/07/2025 \\
			Gilles Morvan & 71 & Calvados & Normandie & 26/09/2026 \\
			Hajar Naccache & 56 & Orne & Normandie & 10/09/2026 \\
			Ibtissem Naccache & 32 & Orne & Normandie & 10/09/2026 \\
			\hline
		\end{tabular}
	}
	\end{table}

	\begin{center}

	\end{center}

\end{frame}

\begin{frame}{Normalisation de la base de client$\cdot$es}

	\begin{figure}[htb]
		\centering
		\includesvg[width=\textwidth]{img/entity-relationship-diagram-example.svg}
	\end{figure}

\end{frame}

\begin{frame}{Normalisation de la base de client$\cdot$es}

	\begin{figure}[htb]
		\resizebox{\textwidth}{!}{
			\begin{tikzpicture}

				\node[label=below:Relation Client$\cdot$es] at (0,0) (client) {
					\begin{tabular}{|c|c|c|c|c|}
						\hline
						\rowcolor{lipstick}\color{white}\textbf{ID} & \color{white}\textbf{Prénom} & \color{white}\textbf{Nom} & \color{white}\textbf{Date de naissance} & \color{white}\textbf{Sexe} \\\hline
						C1000 & Alice & Fleury & 02/08/1996 & F \\
						C1001 & Bob & Brown & 15/03/1987 & M \\
						C1002 & Charlie & O'Brien & 26/10/2000 & F \\
						C1003 & Daniel & Doyle & 11/12/1984 & M \\
						C1004 & Edna & McCarthy & 22/06/1968 & F \\
						C1005 & Fabienne & Lenoir & 09/02/1960 & F \\
						C1006 & Gilles & Morvan & 16/11/1955 & M \\
						C1007 & Hajar & Naccache & 29/05/1970 & F \\
						C1008 & Ibtissem & Naccache & 17/12/1994 & F \\
						\hline
					\end{tabular}
				};

				\node[label=below:Relation Département,right=of client, xshift=0.1cm] (departement) {
					\begin{tabular}{|c|c|c|c|c|c|c|}
						\hline
						\rowcolor{lipstick}\color{white}\textbf{ID} & \color{white}\textbf{Nom département} & \color{white}\textbf{Nom région} \\\hline
						01 & Ain & Auvergne-Rhône-Alpes \\
						02 & Aisne & Hauts-de-France \\
						03 & Allier & Auvergne-Rhône-Alpes \\
						04 & Alpes-de-Haute-Provence & Provence-Alpes-Côte d'Azur \\
						05 & Hautes-Alpes & Provence-Alpes-Côte d'Azur \\
						06 & Alpes-Maritimes & Provence-Alpes-Côte d'Azur \\
						07 & Ardèche & Auvergne-Rhône-Alpes  \\
						08 & ... & ...  \\
						\hline
					\end{tabular}
				};

				\node[label=below:Relation Visite,below=of client, xshift=5cm] (visite) {
					\begin{tabular}{|c|c|c|c|c|}
						\hline
						\rowcolor{lipstick}\color{white}\textbf{ID} & \color{white}\textbf{Date visite} & \color{white}\textbf{ID Client$\cdot$e} &  \color{white}\textbf{ID Département} \\\hline
						V0001 & 01/02/2023 & C1000  & 75 \\
						V0002 & 12/10/2025 & C1001  & 83 \\
						V0002 & 06/11/2026 & C1000  & 83 \\
						V0003 & 17/11/2025 & C1002  & 29 \\
						V0004 & 28/12/2025 & C1003  & 35 \\
						V0005 & 09/01/2026 & C1004  & 60 \\
						V0006 & 26/01/2026 & C1005  & 73 \\
						V0006 & 23/11/2026 & C1005  & 73 \\
						V0007 & 03/02/2026 & C1006  & 14 \\
						V0008 & 10/02/2026 & C1007  & 61 \\
						V0009 & 01/03/2026 & C1008  & 61 \\
						\hline
					\end{tabular}
				};

			\end{tikzpicture}
		}
	\end{figure}

\end{frame}

\begin{frame}{Normalisation et dénormalisation}

    \begin{itemize}
        \item \textbf{Modèle relationnel} : normalisation des données
        \item \textbf{Modèle NoSQL} : dénormalisation des données
    \end{itemize}

    Donner exemple

\end{frame}
